% !TeX root = ../despliegue.tex
% ===========================================================================
%  Variables de entorno.
% ===========================================================================

\section{Configuración}

Todo se configura por variables de entorno, con valores por defecto que
funcionan sin definir ninguna. Las que importan:

\begin{center}
\small
\begin{xltabular}{\textwidth}{@{}l l Y@{}}
\toprule
\textbf{Variable} & \textbf{Por defecto} & \textbf{Qué controla} \\
\midrule
\endfirsthead
\toprule
\textbf{Variable} & \textbf{Por defecto} & \textbf{Qué controla} \\
\midrule
\endhead
\var{POSTGRES_PASSWORD} & \cmd{postgres} &
  Contraseña del superusuario de PostgreSQL \\
\var{USUARIOS_DB_PASSWORD} & \cmd{usuarios_app} &
  Contraseña del rol de \cmd{usuarios_db} \\
\var{CANCHAS_DB_PASSWORD} & \cmd{canchas_app} &
  Contraseña del rol de \cmd{canchas_db} \\
\var{RESERVAS_DB_PASSWORD} & \cmd{reservas_app} &
  Contraseña del rol de \cmd{reservas_db} \\
\var{GATEWAY_PORT} & \cmd{8080} &
  Puerto del gateway en el anfitrión \\
\var{EDGE_PORT} & \cmd{80} &
  Puerto por el que se sirve la aplicación \\
\bottomrule
\end{xltabular}
\end{center}

Se definen en un archivo \cmd{.env} en la raíz del repositorio, que Compose lee
por su cuenta:

\begin{comando}
POSTGRES_PASSWORD=una-contraseña-larga
USUARIOS_DB_PASSWORD=otra
CANCHAS_DB_PASSWORD=otra-más
RESERVAS_DB_PASSWORD=y-otra
\end{comando}

\textbf{Ese archivo no se versiona}: está en \cmd{.gitignore}, de modo que
ninguna credencial llega al repositorio.

\aviso{Las contraseñas de base se aplican \textbf{al crear el volumen}. Si se
cambian con \cmd{pgdata} ya existente, los servicios no podrán conectarse y
reiniciarán en bucle. Hay que hacer \cmd{docker compose down -v} antes.}

Cada microservicio recibe \textbf{solo la credencial de su propia base}.
\servicio{ms-reportes} no recibe ninguna porque no tiene base de datos: obtiene
todo por HTTP de \servicio{ms-reservas} y \servicio{ms-canchas}.
